\documentclass[14pt,a4paper]{extarticle}
\usepackage[T2A]{fontenc}
\usepackage[utf8]{inputenc}
\usepackage[english,russian]{babel}
\usepackage{lmodern}
\usepackage{geometry}
\usepackage{array}
\usepackage{tabularx}
\usepackage{longtable}
\usepackage{ltablex}
\usepackage{booktabs}
\usepackage{hyperref}
\usepackage{indentfirst}
\usepackage{microtype}
\usepackage{graphicx}
\usepackage{setspace}
\usepackage{listings}
\usepackage{graphicx}
\usepackage{caption}
\usepackage{float}
\usepackage{tocloft}
\usepackage{titlesec}
\usepackage{newtxtext,newtxmath}

\titleformat{\section}        % форматируем \section
  {\normalfont\bfseries}      % стиль: обычный + жирный
  {\thesection.}              % номер секции с точкой
  {2mm}                       % расстояние между номером и текстом
  {}                           % перед заголовком можно добавить что-то
\titlespacing*{\section}{1.5cm}{*3}{*1}  % {отступ слева}{отступ сверху}{отступ снизу}

\captionsetup{justification=centering}
\renewcommand{\ttdefault}{cmtt}
\keepXColumns
\setlength{\parindent}{3em}
\newgeometry{left=2.5cm,right=2.5cm,top=3cm,bottom=2.5cm}

\renewcommand{\contentsname}{\hspace*{4em}{\smallСодержание}}
\setlength{\cftbeforesecskip}{0.5em}
\renewcommand{\cftsecleader}{\cftdotfill{\cftdotsep}}
\renewcommand{\cftsecaftersnum}{.\hspace{2mm}}

\title{Отчёт по "Программно-аппаратные средства интеллектуальных систем"}
\author{Выполнил студент группы М-РИТ-25-1, Станиславчук Сергей}

\begin{document}

\setstretch{1.5}
\noindent
\begin{minipage}[t]{0.01\textwidth}
    \includegraphics{~/programmer/++Programmer/5_1/LSTU.jpg}
\end{minipage}%
\begin{minipage}[c]{1.1\linewidth}
    \centering
    {\small
    МИНИСТЕРСТВО НАУКИ И ВЫСШЕГО ОБРАЗОВАНИЯ РФ \\[1em]
    ФЕДЕРАЛЬНОЕ ГОСУДАРСТВЕННОЕ БЮДЖЕТНОЕ \\
    ОБРАЗОВАТЕЛЬНОЕ УЧРЕЖДЕНИЕ ВЫСШЕГО ОБРАЗОВАНИЯ \\
    «ЛИПЕЦКИЙ ГОСУДАРСТВЕННЫЙ ТЕХНИЧЕСКИЙ УНИВЕРСИТЕТ»
    }
\end{minipage}%

\vspace{3em}

\noindent
\hspace{1em}%
\begin{minipage}[t]{0.8\textwidth}
\textbf{Институт} \hspace{2em} компьютерных наук \\
\textbf{Кафедра} \hspace{2em} автоматизированных систем управления
\end{minipage}


\vspace{6em}

\begin{center}
\textbf{Практическая работа №1} \\
по дисциплине \\ <<Планирование карьеры>>
\end{center}

\vspace{6em}

\noindent
\begin{tabular*}{\textwidth}{@{\extracolsep{\fill}} l c r }
    Студент М-РИТ-25-1 & \makebox[4cm]{\hrulefill} & Станиславчук С.М. \\[2em]
Руководитель & \\
Доцент & \makebox[4cm]{\hrulefill} & Триханова И.Е. \\
\end{tabular*}

\vfill
\begin{center}
    Липецк, 2026
\end{center}

\thispagestyle{empty}

\newpage
\tableofcontents
\newpage

\section{Вопрос 3: Можно ли считать смену профессии в 35-40 лет карьерным ростом?}

Смена профессии в возрасте 35–40 лет может рассматриваться как форма карьерного роста, если она обусловлена стремлением человека к профессиональному развитию и самореализации. Карьерный рост не всегда выражается исключительно в повышении должности или заработной платы. В более широком смысле он связан с расширением профессиональных компетенций, ростом мотивации и возможностью реализовать личный потенциал в деятельности, которая вызывает интерес и внутреннюю вовлечённость. Если человек утрачивает интерес к своей прежней профессии, его способность развиваться в данной области неизбежно снижается, поскольку профессиональный рост требует устойчивой мотивации и интеллектуального вовлечения. В такой ситуации переход в другую сферу может стать рациональным и перспективным решением.

Переход в новую профессиональную область способен открыть возможности для развития, которые отсутствовали в прежней деятельности. Интерес к работе формирует основу для повышения квалификации, освоения новых навыков и стремления к достижению более высоких результатов. Следовательно, смена профессии может выступать не как шаг назад, а как стратегическое решение, направленное на долгосрочное развитие карьеры и личностный рост. В данном контексте возраст не является определяющим фактором: ключевое значение имеет способность человека адаптироваться, обучаться и реализовывать свой потенциал в новой сфере.

Показательным примером подобной карьерной трансформации является история японского геймдизайнера Хидетаки Миядзаки. После окончания университета он получил высокооплачиваемую должность персонального менеджера в компании Oracle Corporation. Однако спустя несколько лет работы он пришёл к выводу, что выбранная профессиональная деятельность не приносит ему удовлетворения. Несмотря на стабильность и высокий доход, отсутствие интереса к работе ограничивало его возможности для дальнейшего профессионального развития. Переломным моментом стало знакомство с видеоиграми на консоли PlayStation 3, после чего он начал рассматривать игровые проекты как форму художественного выражения, способную передавать идеи и эмоции автора.

В возрасте около тридцати лет он принял решение оставить престижную должность и перейти в значительно менее известную игровую студию FromSoftware. Такой шаг был особенно рискованным с учётом особенностей японской корпоративной культуры, где традиционно поощряется длительная работа в одной компании. Тем не менее именно этот выбор стал ключевым этапом его профессионального становления. Уже в первые годы работы он занял руководящую позицию в разработке проектов, а затем возглавил работу над игрой Demon's Souls, которая впоследствии стала основой для знаменитой серии Dark Souls. Именно эта серия принесла ему мировую известность и существенно повлияла на развитие современной индустрии видеоигр.

Смена профессии может стать важным этапом карьерного роста. Отказ от стабильной, но неинтересной деятельности в пользу сферы, вызывающей искреннюю заинтересованность, способен привести к значительным профессиональным достижениям.

\newpage
\section{Вопрос 17: Какую роль играет удача в построении карьеры?}

Удача играет заметную, но не определяющую роль в построении карьеры. Под удачей обычно понимают благоприятное стечение обстоятельств: появление подходящей вакансии в нужный момент, знакомство с людьми, способными повлиять на профессиональную траекторию, или участие в проекте, который впоследствии оказывается успешным. Подобные факторы могут ускорить карьерное развитие и предоставить человеку возможность проявить себя в более благоприятных условиях. Однако сами по себе случайные обстоятельства редко становятся достаточным условием для долгосрочного профессионального успеха.

Чтобы воспользоваться удачной возможностью, человек должен обладать определёнными знаниями, интересом к своей деятельности и готовностью действовать. Без этих качеств даже благоприятные обстоятельства не приводят к значительным результатам. Поэтому удача в профессиональной сфере чаще рассматривается как фактор, который усиливает уже существующие способности и стремление человека, а не как основной источник карьерных достижений.

Показательным примером влияния случайных обстоятельств на карьеру является история Геба Ньюэлла — будущего руководителя компании Valve Corporation. Во время академического отпуска после окончания второго курса обучения в Гарвардском Университете он нередко посещал офис компании Microsoft, где работал его брат. Наблюдая за рабочим процессом и общаясь с сотрудниками, Ньюэлл проявлял искренний интерес к деятельности компании и к сфере разработки программного обеспечения. Сотрудники обратили внимание на его заинтересованность и в конечном итоге предложили ему работу. Таким образом, случайное посещение рабочего места родственника стало отправной точкой его профессионального пути в технологической индустрии.

Впоследствии Ньюэлл проработал в Microsoft несколько лет, получив значительный опыт в разработке программных продуктов. Позднее он стал сооснователем компании Valve, которая сыграла важную роль в развитии индустрии видеоигр и цифровой дистрибуции. Этот пример показывает, что удачное стечение обстоятельств может открыть человеку новые возможности, однако решающим фактором остаются личная заинтересованность, готовность учиться и способность использовать появившийся шанс для дальнейшего профессионального развития.

\newpage
\section{Вопрос 19: Почему не каждый хороший специалист становится хорошим руководителем?}

Не каждый высококвалифицированный специалист становится эффективным руководителем, поскольку профессиональная компетентность и управленческие навыки относятся к различным типам деятельности. Специалист обычно сосредоточен на глубоком понимании предметной области, решении технических задач и совершенствовании собственных профессиональных навыков. Руководитель же выполняет иные функции: координирует работу команды, принимает организационные решения, распределяет задачи, управляет конфликтами и отвечает за стратегическое развитие проекта или организации. Поэтому человек может быть выдающимся экспертом в своей сфере, но при этом не обладать склонностью или интересом к управлению людьми.

Кроме того, руководящая деятельность требует развитых коммуникативных навыков, способности делегировать задачи и принимать решения, учитывающие интересы всей команды. Некоторые специалисты предпочитают сохранять концентрацию на технической стороне работы и избегать административных обязанностей, поскольку управление коллективом часто связано с большим количеством организационных и социальных задач. В таких случаях профессионал может осознанно не стремиться к руководящей роли, даже если его компетенции высоко оцениваются в профессиональной среде.

Характерным примером является история разработчика ядра операционной системы Линус Торвальдс и проекта Linux. По мере развития проекта количество разработчиков и объём кода значительно увеличивались, и стало очевидно, что один человек физически не способен контролировать все изменения и поддерживать систему в одиночку. Вместо того чтобы превращать проект в традиционную компанию с офисом, строгой иерархией и постоянным управлением сотрудниками, Торвальдс выбрал иной путь организации работы.

Для координации большого числа разработчиков он создал систему контроля версий Git. Эта технология позволила распределённой группе программистов работать над проектом независимо друг от друга, фиксировать изменения в коде и объединять результаты работы без необходимости централизованного ежедневного управления. Таким образом, вместо классической управленческой модели с жёсткой иерархией была создана технологическая инфраструктура, которая позволила эффективно организовать совместную работу без прямого взаимодействия между людьми (по итогу Линус - специалист в своей области, все еще остается главным разработчиком ядра, над проектом все еще активно работают две тысячи программистов. Линус Торвальдс лишь получает предложения о модификации кода проекта, принимать их или нет - решение остается за ним).

Выдающийся специалист может решать организационные проблемы не через традиционное руководство людьми, а через создание инструментов и систем, упрощающих коллективную работу.
\end{document}
